% Separate three-row layout candidate, 16 November 2026.
% Not installed: all approved figures and the manuscript are preserved.
% PNG files are embedded unchanged; no generative image model was used.
\documentclass[tikz,border=2mm]{standalone}
\usepackage{graphicx}
\usepackage[T1]{fontenc}
\usepackage{helvet}
\usepackage{amsmath,amssymb}
\usepackage{newtxmath}
\usepackage{bm}
\renewcommand{\familydefault}{\sfdefault}
\usetikzlibrary{arrows.meta,calc,decorations.pathmorphing}

\definecolor{ink}{HTML}{202326}
\definecolor{muted}{HTML}{65696D}
\definecolor{hair}{HTML}{B6B8BA}
\definecolor{steel}{HTML}{787C7F}
\definecolor{cyan}{HTML}{477B8B}

% main.tex places the 180 mm figure at 372 TeX pt (130.74 mm).
% 11 TeX pt here gives approximately 8 pt lettering at that final width.
\newcommand{\bodyfont}{\fontsize{11}{12.7}\selectfont}
\newcommand{\smallfont}{\bodyfont}
\newcommand{\headfont}{\bodyfont\bfseries}
\tikzset{
  every node/.style={font=\bodyfont,text=ink,inner sep=0pt},
  flow/.style={draw=ink,line width=.55pt,-{Stealth[length=1.5mm,width=1.0mm]}},
  learning flow/.style={draw=ink,line width=.85pt,-{Stealth[length=1.9mm,width=1.4mm]}},
  spring/.style={draw=steel,line width=.65pt,decorate,
    decoration={zigzag,segment length=1.0mm,amplitude=.42mm,pre length=.55mm,post length=.55mm}},
  guide/.style={draw=hair,line width=.45pt},
}

% Every node represents the same continuous oscillator coordinate.
\newcommand{\oscillator}[2]{%
  \draw[ink,fill=white,line width=.6pt] (#1,#2) circle[radius=1.95mm];
  \fill[cyan] (#1,#2) circle[radius=.65mm];
}

% The force readout and parameter update use the same physical unit.
\newcommand{\oscillatorpair}[2]{%
  \begin{scope}[shift={(#1,#2)}]
    \draw[spring,line width=.85pt] (-3.1,0)--(3.1,0);
    \oscillator{-5}{0}
    \oscillator{5}{0}
  \end{scope}}

\newcommand{\scheduleicon}[2]{%
  \begin{scope}[shift={(#1,#2)}]
    \draw[hair,line width=.4pt] (0,-2.5)--(0,3.5)--(11,3.5);
    \draw[cyan,line width=.85pt,line cap=round]
      (.7,2.7) .. controls (2,-.8) and (6,1.3) .. (8,.4)
      .. controls (9,-.1) and (9.5,-1.7) .. (10.4,-2.0);
    \draw[muted,line width=.75pt,line cap=round]
      (.7,-2.0) .. controls (2.0,.3) and (5.7,-.5) .. (7.2,.5)
      .. controls (8.7,1.0) and (9.4,2.3) .. (10.4,2.8);
  \end{scope}}

\newcommand{\panelheading}[3]{%
  \node[anchor=west,font=\headfont] at (#1,#2) {#3};}

% Miniature full network used for state preparation and short release.
\newcommand{\smallnetwork}[2]{%
  \begin{scope}[shift={(#1,#2)}]
    \tikzset{small spring/.style={draw=steel,line width=.35pt,decorate,
      decoration={zigzag,segment length=.65mm,amplitude=.14mm,
      pre length=.2mm,post length=.2mm}}}
    \foreach \r in {0,1,2} {
      \foreach \c in {0,1,2} {
        \pgfmathsetmacro{\px}{4*\c-.8*\r}
        \pgfmathsetmacro{\py}{2.5*\r}
        \draw[small spring] (\px+.7,\py)--(\px+3.3,\py);
      }
    }
    \foreach \r in {0,1} {
      \foreach \c in {0,1,2,3} {
        \pgfmathsetmacro{\px}{4*\c-.8*\r}
        \pgfmathsetmacro{\py}{2.5*\r}
        \draw[small spring] (\px-.22,\py+.65)--(\px-.58,\py+1.85);
      }
    }
    \foreach \r in {0,1,2} {
      \foreach \c in {0,1,2,3} {
        \pgfmathsetmacro{\px}{4*\c-.8*\r}
        \pgfmathsetmacro{\py}{2.5*\r}
        \begin{scope}[shift={(\px,\py)},scale=.34]
          \oscillator{0}{0}
        \end{scope}
      }
    }
  \end{scope}}

\begin{document}
\begin{tikzpicture}[x=1mm,y=-1mm]
  \path[use as bounding box] (0,0) rectangle (176,160);

  % A. Full-width row: physical network and its symbol legend.
  \panelheading{0}{2.5}{a\quad\mdseries Physical network}
  \draw[muted,line width=.5pt] (5,10) rectangle (123,38.5);
  \foreach \xshift/\idx/\stime in {20/0/\tau,79/1/0} {
    \begin{scope}[shift={(\xshift,16)}]
      \tikzset{network spring/.style={draw=steel,line width=\wire pt,decorate,
        decoration={zigzag,segment length=.9mm,amplitude=\coil mm,
        pre length=.4mm,post length=.4mm}}}
      \foreach \r/\c/\weight in {
        0/0/11,0/1/4,0/2/15,1/0/3,1/1/12,1/2/2,2/0/9,2/1/6,2/2/14} {
        \pgfmathsetmacro{\wire}{(1-\idx)*.45+\idx*(.40+.045*\weight)}
        \pgfmathsetmacro{\coil}{(1-\idx)*.22+\idx*(.19+.024*\weight)}
        \pgfmathsetmacro{\px}{3+9*\c-1.5*\r}
        \pgfmathsetmacro{\py}{5.8*\r}
        \draw[network spring] (\px+1.6,\py)--(\px+7.4,\py);
      }
      \foreach \r/\c/\weight in {
        0/0/16,0/1/5,0/2/8,0/3/7,1/0/10,1/1/0,1/2/13,1/3/1} {
        \pgfmathsetmacro{\wire}{(1-\idx)*.45+\idx*(.40+.045*\weight)}
        \pgfmathsetmacro{\coil}{(1-\idx)*.22+\idx*(.19+.024*\weight)}
        \pgfmathsetmacro{\px}{3+9*\c-1.5*\r}
        \pgfmathsetmacro{\py}{5.8*\r}
        \draw[network spring] (\px-.4,\py+1.45)--(\px-1.1,\py+4.35);
      }
      \foreach \r in {0,1,2} {
        \foreach \c in {0,1,2,3} {
          \pgfmathsetmacro{\px}{3+9*\c-1.5*\r}
          \pgfmathsetmacro{\py}{5.8*\r}
          \begin{scope}[shift={(\px,\py)},scale=.8]
            \oscillator{0}{0}
          \end{scope}
        }
      }
      \node[text=muted] at (15,17.5) {$s=\stime$};
    \end{scope}
  }
  \node at (10,15) {$T$};
  \node[text=muted] at (64,20) {$\theta(\tau-s)$};
  \node[text=muted] at (64,26) {$\cdots$};
  \node[text=muted,fill=white,inner xsep=1mm] at (64,38.5) {Thermal bath};
  \begin{scope}[shift={(131.5,18)},scale=.8]
    \oscillator{0}{0}
  \end{scope}
  \node[anchor=west] at (137,18) {Oscillators};
  \draw[spring,line width=.45pt] (128,29)--(135,29);
  \node[anchor=west,align=left] at (138,29) {Time evolving\\couplings};

  % B. Four independent paths connect the common clamp and update stages.
  \panelheading{0}{47}{b\quad\mdseries Local learning}

  % The full state is prepared, with a real snapshot reused as a visual example.
  \begin{scope}[shift={(8.5,74)},xslant=.25]
    \node[anchor=north west,transform shape] at (0,0)
      {\includegraphics[width=23.1mm,height=12.6mm]{images/frame_traj00_s4.25_autoscale.png}};
  \end{scope}
  \foreach \x in {8.3,14.3,26.3} {
    \draw[white,line width=2pt] (\x,82.8)--(\x,70.8);
    \draw[cyan,line width=1.05pt,-{Stealth[length=1.8mm,width=1.3mm]}]
      (\x,82.8)--(\x,70.8);
  }
  \begin{scope}[shift={(18.5,65)},scale=1.5]
    \smallnetwork{-5.2}{-2.5}
  \end{scope}
  \node[anchor=north,align=center] at (20,91) {Clamp\\full state};
  \node[anchor=north,text=muted] at (20,103) {$\bm x_t\sim p_t$};

  \draw[muted,line width=.65pt] (51,54.5) rectangle (84,89.5);
  \draw[muted,line width=.65pt] (99,54.5) rectangle (132,89.5);
  \node[text=muted] at (91.5,79) {or};

  \node at (67.5,59) {Force matching};
  \oscillatorpair{67.5}{70.5}
  \foreach \x/\idx in {62.5/n,72.5/m} {
    \draw[ink,line width=.65pt,line cap=round] (\x-2.7,67.5)--(\x+2.7,67.5);
    \draw[ink,line width=.65pt,line cap=round] (\x-2.7,73.5)--(\x+2.7,73.5);
    \node[text=cyan] at (\x,64.2) {$\widehat F_{\idx}$};
  }
  \node[anchor=north,align=center] at (67.5,78) {Read local\\forces};

  \node at (115.5,59) {CD1};
  \begin{scope}[shift={(115.5,70.5)},scale=1.5]
    \smallnetwork{-5.2}{-2.5}
    \draw[flow,draw=cyan] (-5.4,-3.5)--(-6.9,-5.1);
    \draw[flow,draw=cyan] (7.4,-2.8)--(9,-4.5);
    \draw[flow,draw=cyan] (1.2,3.6)--(3.1,5);
  \end{scope}
  \node[anchor=north] at (115.5,81) {Release network};

  % Direct path to force readout, routed path to CD1; neither crosses a box.
  \draw[learning flow] (36,77)--(50,72);
  \draw[learning flow] (36,79)--(43,97)--(115.5,97)--(115.5,90.5);
  \draw[learning flow] (85,70.5)--(89,51)--(146,51)--(151,68.5);
  \draw[learning flow] (133,74)--(151,72.5);

  \oscillatorpair{160}{70.5}
  \node[text=cyan] at (160,64) {$\theta_i$};
  \draw[flow,draw=cyan] (160,65.8)--(160,68.3);
  \node[anchor=north,align=center] at (160,80) {Update\\parameters};
  \draw[learning flow] (160,92)--(160,95.5);
  \scheduleicon{154.5}{100}
  \node[anchor=north] at (158,106) {Learned schedule};
  \node[anchor=north] at (158,111) {$\bm\theta^*(t)$};

  % C. Author-supplied snapshots retain their original pixel data and order.
  \panelheading{0}{119}{c\quad\mdseries Generate}
  \scheduleicon{59}{118.5}
  \node[anchor=west] at (75,119) {External replay $\bm\theta^*(\tau-s)$};
  \draw[flow,draw=muted] (16,125)--(167,125);
  \node[anchor=west] at (170,125) {$s$};
  \foreach \stime/\cx in {0.00/16,3.25/52,3.75/88,4.25/124,5.00/160} {
    \node[anchor=south,text=muted] at (\cx,131) {$s=\stime$};
    \node[anchor=north west] at (\cx-10.5,133)
      {\includegraphics[width=21mm]{images/frame_traj00_s\stime_autoscale.png}};
  }
  \node[anchor=north] at (16,157) {Noisy state};
  \node[anchor=north] at (160,157) {Sample};
\end{tikzpicture}
\end{document}
